\scanpage{88}\mspage{T-2}

On pose également:
\[
 c_{4}=b_{2}^{2}-24\,b_{4}\qquad -c_{6}=b_{2}^{3}-36b_{2}b_{4}+216b_{6}
\]

(Les $c_{i}$ sont liés aux classiques $g_{2}$ et $g_{3}$ par les formule\rlap{\hspace{-.16em}\handwritten{\tikz[baseline=0pt]{\draw[line width=.32pt] (0,.04ex) .. controls (.18em,-.30ex) and (.45em,-.52ex) .. (.73em,-.71ex);}}}s:
\[
 c_{4}=\underset{\handwritten{2^{\unsure{2'|2}}\mathord{.}3}}{12}\,g_{2}
 \qquad c_{6}=\underset{\handwritten{2^{3}\mathord{.}3^{3}}}{216}\,g_{3}\ .)
\]

Et :
\[
\begin{aligned}
\handwritten{\Delta}&=-b_{2}^{2}b_{8}-8b_{4}^{3}-27b_{6}^{2}+9b_{2}b_{4}b_{6}\\
&=b_{4}^{3}-27b_{6}^{2}+b_{8}(36b_{4}-b_{2}^{2}).
\end{aligned}
\]

On a :
\[
 \underset{\handwritten{2^{6}\,3^{3}}}{1728}\,\handwritten{\Delta}=c_{4}^{3}-c_{6}^{2}\ .
\]

On pose enfin (pour $\handwritten{\Delta}$ inversible) :
\[
 j=\frac{c_{4}^{3}}{\handwritten{\Delta}}\qquad
 (j-1728)\handwritten{\Delta}=c_{6}^{2}\qquad
 \frac{j}{j-1728}=\frac{c_{4}^{3}}{c_{6}^{2}}
\]

Si l'on pose $F=y^{2}+a_{1}xy+a_{3}y-x^{3}-a_{2}x^{2}-a_{4}x-a_{6}$, on a :
\[
 F_{y}=2y+a_{1}x+a_{3}\qquad -F_{x}=3x^{2}+2a_{2}x+a_{4}-a_{1}y\ .
\]

Si $(\ast)$ définit bien une courbe elliptique (i.e. si $\handwritten{\Delta}$ est inversible), on peut écrire la forme différentielle $\pi$ de départ sous la forme $\pi=dx/F_{y}=-dy/F_{x}$, i.e. :
\[
 \pi=\frac{dx}{2y+a_{1}x+a_{3}}=\frac{dy}{3x^{2}+2a_{2}x+a_{4}-a_{1}y}
\]
% archive sheet number: 87
